% ============================================================
% 运筹学优化算法 — Beamer 课件
% Topics: 线性规划 · 非线性规划 · 凸优化 · ML / LLM 应用
% Compile: xelatex optimization_beamer.tex
% ============================================================
\documentclass[11pt, aspectratio=169, svgnames, table, fontset=fandol]{ctexbeamer}

% ---------- Packages ----------
\usepackage{amsmath, amssymb, amsthm}
\usepackage{mathrsfs}
\usepackage{bm}
\usepackage{booktabs}
\usepackage{multirow}
\usepackage{tikz}
\usetikzlibrary{arrows.meta, positioning, calc, shapes.geometric}
\usepackage{algorithm}
\usepackage{algpseudocode}
\usepackage{graphicx}

% ---------- Theme ----------
\usetheme{Madrid}
\setbeamertemplate{frametitle continuation}{}
% --- Block colors: light bg + dark text + border (no black rendering risk) ---
\setbeamercolor{block title}{fg=blue!40!black,bg=blue!10!white}
\setbeamercolor{block body}{fg=black,bg=blue!3!white}
\setbeamercolor{block title alerted}{fg=red!40!black,bg=red!10!white}
\setbeamercolor{block body alerted}{fg=black,bg=red!3!white}
\setbeamercolor{block title example}{fg=green!40!black,bg=green!10!white}
\setbeamercolor{block body example}{fg=black,bg=green!3!white}
\setbeamertemplate{blocks}[default]

\setbeamertemplate{navigation symbols}{}
\setbeamertemplate{footline}[frame number]
\setbeamertemplate{headline}{}

% ---------- Custom ----------
\setbeamercolor{alerted text}{fg=red!70!black}
\newcommand{\R}{\mathbb{R}}
\newcommand{\E}{\mathbb{E}}
\newcommand{\argmin}{\operatorname*{arg\,min}}
\newcommand{\argmax}{\operatorname*{arg\,max}}
\newcommand{\diag}{\operatorname{diag}}
\newcommand{\tr}{\operatorname{tr}}
\newcommand{\prox}{\operatorname{prox}}
\newcommand{\dom}{\operatorname{dom}}
\newcommand{\epi}{\operatorname{epi}}

% ---------- Title ----------
\title{\textbf{运筹学优化算法}}
\subtitle{线性规划 · 非线性规划 · 凸优化 · 机器学习与大模型应用}
\author{运筹学课件}
\institute{Operations Research \& Optimization}
\date{\today}

\begin{document}

% ============================================================
% TITLE SLIDE
% ============================================================
\begin{frame}
  \titlepage
\end{frame}

% ============================================================
% OUTLINE
% ============================================================
\begin{frame}{目录}
  \begin{description}
    \item[Part 1] 线性规划 (Linear Programming)
    \item[Part 2] 非线性规划 (Nonlinear Programming)
    \item[Part 2.5] 启发式与元启发式算法 (Heuristics \& Metaheuristics)
    \item[Part 3] 凸优化 (Convex Optimization)
    \item[Part 4] 运筹优化在 ML 与大模型中的应用
    \item[Part 5] GPU / CUDA 并行计算与 ML 优化加速
  \end{description}

  \vspace{6pt}
  \begin{center}
    \small 共计 61 张幻灯片
  \end{center}
\end{frame}

% ############################################################
% PART 1: 线性规划
% ############################################################
\section{线性规划 (Linear Programming)}

\subsection{LP 基本概念与标准形}

\begin{frame}{线性规划简介}
  \textbf{线性规划 (LP)}是运筹学中最基础也是最重要的优化模型。
  \begin{itemize}
    \item 目标函数是决策变量的线性函数
    \item 约束条件是一组线性等式/不等式
    \item 决策变量可以取连续值（早期经典 LP 中为非负）
  \end{itemize}

  \vspace{6pt}
  \textbf{标准形 (Standard Form)}：
  \begin{align*}
    \min_{\bm{x}} \quad & \bm{c}^\top \bm{x} \\
    \text{s.t.} \quad  & \bm{A} \bm{x} = \bm{b} \\
                       & \bm{x} \geq \bm{0}
  \end{align*}
  其中 $\bm{A} \in \R^{m \times n}$，$\bm{b} \in \R^m$，$\bm{c} \in \R^n$。
\end{frame}

\begin{frame}{LP 的几何意义}
  \begin{itemize}
    \item 可行域是\textbf{多面体 (Polyhedron)}：$\{\bm{x} \mid \bm{Ax} = \bm{b},\ \bm{x} \geq \bm{0}\}$
    \item 若可行域非空且有界，则为凸多胞形 (Convex Polytope)
    \item \textbf{最优解必在顶点 (extreme point / basic feasible solution) 处取得}
    \item 单纯形法 (Simplex Method) 正是沿着顶点间迭代
  \end{itemize}
  \vspace{10pt}
  \begin{center}
    \begin{tikzpicture}[scale=0.75]
      \draw[->] (-0.5,0) -- (5.5,0) node[right] {$x_1$};
      \draw[->] (0,-0.5) -- (0,4.5) node[above] {$x_2$};
      \filldraw[blue!15, draw=blue!50] (0,0) -- (4,0) -- (3,3) -- (0,3) -- cycle;
      \draw[->, red, thick] (0.2,0.2) -- (1.5,2.5);
      \node[red] at (1.8,2.8) {$\bm{c}$};
      \node at (1.5,1.5) {可行域};
      \node[circle, fill=red, inner sep=2pt] at (3,3) {};
      \node[above] at (3,3) {最优解};
    \end{tikzpicture}
  \end{center}
\end{frame}

\begin{frame}{LP 的变形}
  任何 LP 都可以转化为标准形：
  \begin{itemize}
    \item $\max \bm{c}^\top \bm{x} \;\Longleftrightarrow\; \min -\bm{c}^\top \bm{x}$
    \item $\bm{Ax} \leq \bm{b} \;\Longleftrightarrow\; \bm{Ax} + \bm{s} = \bm{b},\ \bm{s} \geq \bm{0}$（松弛变量）
    \item $\bm{Ax} \geq \bm{b} \;\Longleftrightarrow\; \bm{Ax} - \bm{s} = \bm{b},\ \bm{s} \geq \bm{0}$（剩余变量）
    \item $x_i$ 无约束 $\;\Longleftrightarrow\; x_i = x_i^+ - x_i^-,\ x_i^+, x_i^- \geq 0$
  \end{itemize}
  \vspace{8pt}
  \begin{alertblock}{重要性质}
    LP 可在多项式时间内求解（椭球法 / 内点法），但 Simplex 法在实践中仍然极其高效。
  \end{alertblock}
\end{frame}

\subsection{单纯形法 (Simplex Method)}

\begin{frame}{单纯形法基本思路}
  \textbf{基本可行解 (Basic Feasible Solution, BFS)}：$\bm{A}$ 的列中选 $m$ 个线性无关列构成基矩阵 $\bm{B}$，令非基变量为 0，解 $\bm{Bx}_B = \bm{b}$。

  \vspace{8pt}
  \textbf{单纯形法迭代步骤}：
  \begin{enumerate}
    \item 从一个初始 BFS 出发
    \item 计算简约代价 (reduced cost)：$\bar{\bm{c}}_N^\top = \bm{c}_N^\top - \bm{c}_B^\top \bm{B}^{-1}\bm{N}$
    \item 若 $\bar{\bm{c}}_N \geq \bm{0}$（极小化问题），当前解最优
    \item 否则选入基变量（Dantzig 规则：选最负），出基变量经最小比值 (ratio test) 确定
    \item 更新基矩阵，进入下一轮迭代
  \end{enumerate}
  几何上等价于沿着多面体的棱在顶点间移动。
\end{frame}

\begin{frame}{单纯形表 (Simplex Tableau)}
  标准形 LP 的初始单纯形表：
  \begin{center}
    \begin{tabular}{c|cc|c}
      \toprule
      基变量 & $\bm{x}_B$ & $\bm{x}_N$ & RHS \\
      \midrule
      $\bm{x}_B$ & $\bm{I}$ & $\bm{B}^{-1}\bm{N}$ & $\bm{B}^{-1}\bm{b}$ \\
      \midrule
      $z$ & $\bm{0}$ & $\bar{\bm{c}}_N^\top$ & $\bm{c}_B^\top \bm{B}^{-1}\bm{b}$ \\
      \bottomrule
    \end{tabular}
  \end{center}
  \begin{itemize}
    \item $\bar{\bm{c}}_N \geq \bm{0}$ 时达到最优
    \item 若有 $\bar{c}_j < 0$ 且对应列 $\leq \bm{0}$，则问题无界
    \item Min ratio test 保证可行性
  \end{itemize}
\end{frame}

\begin{frame}{两阶段法 (Two-Phase Method)}
  初始 BFS 通常不易直接获得。\textbf{两阶段法}通过引入人工变量解决：
  \begin{itemize}
    \item \textbf{Phase I}：引入人工变量 $\bm{a}$，求解辅助问题
      \begin{equation*}
        \min \sum_{i} a_i \quad \text{s.t.} \quad \bm{Ax} + \bm{a} = \bm{b},\ \bm{x}, \bm{a} \geq \bm{0}
      \end{equation*}
      若最优值 $>0$，原问题无可行解
    \item \textbf{Phase II}：以 Phase I 的最优 BFS 为起点，求解原问题
    \item 大 M 法是其近似替代（在目标函数中加入惩罚项 $M\sum a_i$）
  \end{itemize}
\end{frame}

\subsection{对偶理论 (Duality Theory)}

\begin{frame}{对偶理论：构造对偶问题}
  原始问题 (Primal)：
  \begin{equation*}
    \min_{\bm{x}} \;\bm{c}^\top \bm{x} \quad \text{s.t.} \quad \bm{Ax} = \bm{b},\ \bm{x} \geq \bm{0}
  \end{equation*}

  \vspace{6pt}
  对偶问题 (Dual)：
  \begin{equation*}
    \max_{\bm{y}} \;\bm{b}^\top \bm{y} \quad \text{s.t.} \quad \bm{A}^\top \bm{y} \leq \bm{c}
  \end{equation*}

  \vspace{8pt}
  \begin{table}
    \small
    \begin{tabular}{c|c}
      \toprule
      Primal (极小化) & Dual (极大化) \\
      \midrule
      变量 $x_i \geq 0$ & $\bm{a}_i^\top \bm{y} \leq c_i$ \\
      变量 $x_i$ 无约束 & $\bm{a}_i^\top \bm{y} = c_i$ \\
      $\bm{a}_i^\top \bm{x} \geq b_i$ & 变量 $y_i \geq 0$ \\
      $\bm{a}_i^\top \bm{x} = b_i$ & 变量 $y_i$ 无约束 \\
      \bottomrule
    \end{tabular}
  \end{table}
\end{frame}

\begin{frame}{弱对偶与强对偶}
  \begin{block}{弱对偶定理 (Weak Duality)}
    对于任意可行 $\bm{x}$（Primal）和 $\bm{y}$（Dual），有：
    \begin{equation*}
      \bm{c}^\top \bm{x} \geq \bm{b}^\top \bm{y}
    \end{equation*}
    对偶问题最优值给出原始问题最优值的下界。
  \end{block}

  \begin{block}{强对偶定理 (Strong Duality)}
    若 Primal 有有限最优解，则 Dual 也有最优解，且：
    \begin{equation*}
      \min \bm{c}^\top \bm{x} = \max \bm{b}^\top \bm{y}
    \end{equation*}
  \end{block}

  \begin{alertblock}{互补松弛条件 (Complementary Slackness)}
    $x_i^* \cdot (\bm{a}_i^\top \bm{y}^* - c_i) = 0$ 对所有 $i$；\\[2pt]
    $y_i^* \cdot (\bm{a}_i^\top \bm{x}^* - b_i) = 0$ 对所有 $i$。
  \end{alertblock}
\end{frame}

\begin{frame}{对偶理论的经济解释：影子价格}
  \begin{itemize}
    \item 对偶变量 $\bm{y}$ 也称\textbf{影子价格 (Shadow Price)}
    \item $y_i$ = 第 $i$ 个约束右端项 $b_i$ 每增加 1 单位，最优目标值的边际改善
    \item 影子价格为 0 的约束是非紧的（有松弛）
    \item 影子价格 $>0$ 的约束是紧的（资源稀缺）
  \end{itemize}

  \vspace{8pt}
  \textbf{灵敏度分析 (Sensitivity Analysis)} 即利用对偶信息，判断 $b_i$ 或 $c_j$ 在什么范围内波动后当前基仍然最优。

  \begin{itemize}
    \item 右端项范围：$\bm{B}^{-1}\bm{b} \geq \bm{0}$ 的范围
    \item 目标系数范围：$\bar{\bm{c}}_N \geq \bm{0}$ 的范围
  \end{itemize}
\end{frame}

\subsection{内点法 (Interior Point Methods)}

\begin{frame}{线性规划的内点法}
  LP 也可以用\textbf{内点法 (Interior Point Method)}求解，在可行域内部迭代，多项式时间保证。

  \vspace{6pt}
  \textbf{原始-对偶路径跟踪法}：解 KKT 系统的扰动形式
  \begin{align*}
    \bm{Ax} &= \bm{b}, \quad \bm{x} > \bm{0} \\
    \bm{A}^\top \bm{y} + \bm{s} &= \bm{c}, \quad \bm{s} > \bm{0} \\
    \bm{X} \bm{S} \bm{e} &= \mu \bm{e}
  \end{align*}
  其中 $\bm{X} = \diag(\bm{x}),\ \bm{S} = \diag(\bm{s})$，$\mu \to 0$。

  \vspace{6pt}
  用 Newton 法求解该系统，逐步降低 $\mu$。实践中对大规模稀疏 LP 效果极好。
\end{frame}

\begin{frame}{LP 应用举例}
  \begin{itemize}
    \item \textbf{生产计划}：多产品多资源在约束下最大化利润
    \item \textbf{运输问题}：$m$ 个产地、$n$ 个销地的最小成本运输方案
    \item \textbf{分配问题}：Hungarian 算法求解 $n$ 对 $n$ 的最优匹配
    \item \textbf{网络流问题}：最大流 / 最小费用流 → 特殊 LP 结构
    \item \textbf{投资组合}：线性约束下的资产配置
    \item \textbf{数据包络分析 (DEA)}：效率评价
  \end{itemize}
  \vspace{6pt}
  \begin{center}
    \textcolor{blue}{大规模 LP 求解器：Gurobi, CPLEX, HiGHS, GLPK}
  \end{center}
\end{frame}

% ############################################################
% PART 2: 非线性规划
% ############################################################
\section{非线性规划 (Nonlinear Programming)}

\subsection{NLP 基础概念}

\begin{frame}{非线性规划简介}
  NLP 的标准形式：
  \begin{align*}
    \min_{\bm{x}} \quad & f(\bm{x}) \\
    \text{s.t.} \quad   & g_i(\bm{x}) \leq 0,\quad i = 1,\dots,m \\
                        & h_j(\bm{x}) = 0,\quad j = 1,\dots,p
  \end{align*}
  其中 $f, g_i, h_j$ 至少有一个是非线性的。

  \vspace{8pt}
  \begin{itemize}
    \item \textbf{无约束 NLP}：只需 $\min f(\bm{x})$
    \item \textbf{约束 NLP}：需要处理等式和不等式约束
    \item 解可能非全局最优 → 需要全局优化技术
  \end{itemize}
\end{frame}

\subsection{无约束优化}

\begin{frame}{无约束最优性条件}
  对于 $\min f(\bm{x}),\ \bm{x} \in \R^n$：

  \begin{block}{一阶必要条件}
    若 $\bm{x}^*$ 是局部极小点且 $f$ 可微，则 $\nabla f(\bm{x}^*) = \bm{0}$。
  \end{block}

  \begin{block}{二阶必要条件}
    若 $\bm{x}^*$ 是局部极小点且 $f$ 二阶可微，则 $\nabla f(\bm{x}^*) = \bm{0}$ 且 $\nabla^2 f(\bm{x}^*) \succcurlyeq \bm{0}$（半正定）。
  \end{block}

  \begin{block}{二阶充分条件}
    若 $\nabla f(\bm{x}^*) = \bm{0}$ 且 $\nabla^2 f(\bm{x}^*) \succ \bm{0}$（严格正定），则 $\bm{x}^*$ 是严格局部极小点。
  \end{block}
\end{frame}

\begin{frame}{梯度下降法 (Gradient Descent)}
  迭代形式：
  \begin{equation*}
    \bm{x}_{k+1} = \bm{x}_k - \alpha_k \nabla f(\bm{x}_k)
  \end{equation*}

  \begin{itemize}
    \item $\alpha_k$：步长 (step size / learning rate)
    \item 收敛速度：线性收敛（对一般凸函数），$\mathcal{O}(1/k)$ 次梯度收敛
    \item \textbf{Wolfe 条件}确保充分下降和曲率
    \item \textbf{Backtracking line search} 是常用的步长选择策略
  \end{itemize}

  \vspace{6pt}
  \textcolor{red}{缺点}：收敛慢，对 ill-conditioned 问题表现差，需精细调整步长。
\end{frame}

\begin{frame}{牛顿法 (Newton's Method)}
  利用二阶 Taylor 展开：
  \begin{equation*}
    f(\bm{x}_k + \bm{p}) \approx f(\bm{x}_k) + \nabla f(\bm{x}_k)^\top \bm{p} + \frac12 \bm{p}^\top \nabla^2 f(\bm{x}_k) \bm{p}
  \end{equation*}

  Newton 方向：$\bm{p}_k = -[\nabla^2 f(\bm{x}_k)]^{-1} \nabla f(\bm{x}_k)$
  \begin{equation*}
    \bm{x}_{k+1} = \bm{x}_k + \alpha_k \bm{p}_k
  \end{equation*}

  \begin{itemize}
    \item 二次终止性：对严格凸二次函数一步到最优
    \item 局部二次收敛（极快）
    \item \textcolor{red}{缺点}：每次迭代需计算和分解 Hessian（$\mathcal{O}(n^3)$），Hessian 可能非正定
  \end{itemize}
\end{frame}

\begin{frame}{拟牛顿法 (Quasi-Newton Methods)}
  避免计算 Hessian，用梯度信息近似 Hessian 逆。

  \vspace{6pt}
  \textbf{BFGS (Broyden–Fletcher–Goldfarb–Shanno)}：
  \begin{align*}
    s_k &= \bm{x}_{k+1} - \bm{x}_k \\
    y_k &= \nabla f(\bm{x}_{k+1}) - \nabla f(\bm{x}_k) \\
    \bm{H}_{k+1} &= \left(\bm{I} - \rho_k s_k y_k^\top\right) \bm{H}_k \left(\bm{I} - \rho_k y_k s_k^\top\right) + \rho_k s_k s_k^\top
  \end{align*}
  其中 $\rho_k = 1 / (y_k^\top s_k)$。

  \begin{itemize}
    \item 超线性收敛
    \item \textbf{L-BFGS}：仅存储最近 $m$ 对 $(s_k, y_k)$，内存 $\mathcal{O}(mn)$，适合大规模问题
  \end{itemize}
\end{frame}

\subsection{约束优化与 KKT 条件}

\begin{frame}{Karush–Kuhn–Tucker (KKT) 条件}
  NLP 的一阶必要条件——KKT 条件：
  \begin{align*}
    \nabla f(\bm{x}^*) + \sum_{i=1}^m \mu_i \nabla g_i(\bm{x}^*) + \sum_{j=1}^p \lambda_j \nabla h_j(\bm{x}^*) &= \bm{0} \\
    g_i(\bm{x}^*) &\leq 0,\quad i=1,\dots,m \\
    h_j(\bm{x}^*) &= 0,\quad j=1,\dots,p \\
    \mu_i &\geq 0,\quad i=1,\dots,m \\
    \mu_i \cdot g_i(\bm{x}^*) &= 0,\quad i=1,\dots,m \quad \text{(互补松弛)}
  \end{align*}

  \begin{itemize}
    \item $\mu_i$：不等式约束的 Lagrange 乘子
    \item $\lambda_j$：等式约束的 Lagrange 乘子
    \item 在凸规划中，KKT 条件既是必要也是充分的
  \end{itemize}
\end{frame}

\begin{frame}{约束优化的数值方法}
  \begin{itemize}
    \item \textbf{罚函数法 (Penalty Method)}：将约束惩罚加入目标
      \begin{equation*}
        \min_{\bm{x}} f(\bm{x}) + \rho \left[\sum_i \max(0, g_i(\bm{x}))^2 + \sum_j h_j(\bm{x})^2\right]
      \end{equation*}
      $\rho \to \infty$ 时收敛，但可能导致 ill-conditioning

    \item \textbf{障碍函数法 (Barrier Method)}：用对数障碍处理不等式约束
      \begin{equation*}
        \min_{\bm{x}} f(\bm{x}) - \mu \sum_i \ln(-g_i(\bm{x}))
      \end{equation*}
      $\mu \to 0$ 时收敛

    \item \textbf{SQP (Sequential Quadratic Programming)}：每步解一个二次规划 (QP) 子问题
    \item \textbf{增广 Lagrange 法 (Augmented Lagrangian)}：结合罚函数和对偶思想
  \end{itemize}
\end{frame}

\begin{frame}{信赖域方法 (Trust Region Methods)}
  每步在一个信赖域半径 $\Delta_k$ 内近似目标函数：
  \begin{equation*}
    \min_{\|\bm{p}\| \leq \Delta_k} m_k(\bm{p}) = f(\bm{x}_k) + \nabla f(\bm{x}_k)^\top \bm{p} + \frac12 \bm{p}^\top \bm{B}_k \bm{p}
  \end{equation*}

  \vspace{4pt}
  根据实际下降与模型预测下降的比率 $\rho_k$ 决定是否接受步及如何调整 $\Delta_k$：
  \begin{equation*}
    \rho_k = \frac{f(\bm{x}_k) - f(\bm{x}_k + \bm{p}_k)}{m_k(\bm{0}) - m_k(\bm{p}_k)}
  \end{equation*}

  \begin{itemize}
    \item 比线搜索更稳健，能自动调节步长
    \item 适合非凸和 ill-conditioned 问题
    \item 在大规模机器学习中也有应用（信赖域 Newton 法）
  \end{itemize}
\end{frame}

% ############################################################

% ############################################################
% PART 2.5: 启发式与元启发式算法
% ############################################################
\section{启发式与元启发式算法 (Heuristics \& Metaheuristics)}

\subsection{简介与分类}

\begin{frame}{为什么需要启发式算法？}
  NLP 中许多实际问题是非凸、NP-hard 的，精确算法在可接受时间内无法求解。

  \vspace{6pt}
  \textbf{启发式 (Heuristic)}：基于问题特定知识的经验规则，能快速找到较好解，但不保证最优。

  \vspace{4pt}
  \textbf{元启发式 (Metaheuristic)}：问题无关的高层搜索策略，指导底层启发式探索解空间。

  \vspace{6pt}
  \begin{block}{核心思想}
    在 \textbf{exploitation (利用)} 和 \textbf{exploration (探索)} 之间平衡，避免陷入局部最优。
  \end{block}

  \begin{columns}[T]
    \column{0.48\textwidth}
    \textbf{单点搜索}：
    \begin{itemize}
      \item 局部搜索 / 爬山法
      \item 模拟退火 (SA)
      \item 禁忌搜索 (TS)
      \item 变邻域搜索 (VNS)
    \end{itemize}
    \column{0.48\textwidth}
    \textbf{群体搜索}：
    \begin{itemize}
      \item 遗传算法 (GA)
      \item 粒子群优化 (PSO)
      \item 蚁群算法 (ACO)
      \item 差分进化 (DE)
    \end{itemize}
  \end{columns}
\end{frame}

\begin{frame}{贪心算法与局部搜索}
  \textbf{贪心算法 (Greedy)}：每一步选择当前最优，不回溯。速度快但易入局部最优。

  \vspace{4pt}
  \textbf{局部搜索 (Local Search) / 爬山法 (Hill Climbing)}：
  \begin{itemize}
    \item 从初始解出发，在邻域 $N(\bm{x})$ 中迭代改进
    \item 若 $f(\bm{x}') < f(\bm{x})$，接受 $\bm{x}'$，否则停止
    \item 典型结果：局部最优 —— 无法通过邻域移动找到更好解
  \end{itemize}

  \vspace{6pt}
  \begin{alertblock}{困境}
    纯局部搜索在复杂的“崎岖”解空间中极易卡在低质量局部最优。元启发式通过接受“坏”移动来逃脱局部最优。
  \end{alertblock}
\end{frame}

\subsection{模拟退火 (Simulated Annealing)}

\begin{frame}{模拟退火 (SA)}
  模拟金属退火的物理过程：高温 → 原子活跃跳跃 → 缓慢冷却 → 结晶到最低能态。

  \vspace{4pt}
  \textbf{Metropolis 准则}（概率接受劣解）：
  \begin{itemize}
    \item 若 $\Delta E = f(\bm{x}_{\text{new}}) - f(\bm{x}) < 0$：接受新解
    \item 若 $\Delta E \geq 0$：以概率 $p = \exp(-\Delta E / T)$ 接受
  \end{itemize}
  温度 $T$ 从高逐渐降低（冷却调度 cooling schedule），早期 $T$ 高时几乎接受所有移动。

  \vspace{6pt}
  \textbf{关键参数}：
  \begin{itemize}
    \item 初始温度 $T_0$：需足够高
    \item 冷却率 $\alpha$：$T_{k+1} = \alpha \cdot T_k,\ \alpha \in (0.9, 0.99)$
    \item 终止条件：$T < T_{\min}$ 或最大迭代次数
  \end{itemize}

  理论上以概率 1 收敛到全局最优（当冷却足够缓慢），但实践中需要合理调参。
\end{frame}

\subsection{遗传算法 (Genetic Algorithm)}

\begin{frame}{遗传算法 (GA)}
  模拟自然选择和遗传进化过程。

  \vspace{4pt}
  \textbf{基本流程}：
  \begin{enumerate}
    \item \textbf{初始化}：随机生成种群 (population)
    \item \textbf{适应度评估}：计算每个个体的 $f(\bm{x})$
    \item \textbf{选择 (Selection)}：轮盘赌 / 锦标赛选择 —— 优者更可能被选中
    \item \textbf{交叉 (Crossover)}：两个父代交换基因片段产生子代
    \item \textbf{变异 (Mutation)}：以低概率随机扰动基因，维持多样性
    \item \textbf{替换}：新一代取代旧一代，重复 2--5
  \end{enumerate}

  \vspace{4pt}
  \textbf{编码方式}：二进制编码、实数编码、排列编码（TSP 等组合问题）。

  \begin{itemize}
    \item 交叉率 $p_c \approx 0.7\text{--}0.9$，变异率 $p_m \approx 0.001\text{--}0.1$
    \item 精英策略 (Elitism)：保留最优个体直接进入下一代
  \end{itemize}
\end{frame}

\subsection{粒子群优化 (PSO)}

\begin{frame}{粒子群优化 (PSO)}
  模拟鸟群、鱼群的群体协作行为。

  \vspace{4pt}
  每个粒子 $i$ 在搜索空间中飞行，记录：
  \begin{itemize}
    \item $\bm{x}_i$：当前位置
    \item $\bm{v}_i$：速度
    \item $\bm{p}_{\text{best},i}$：个人历史最优位置
    \item $\bm{g}_{\text{best}}$：全局最优位置
  \end{itemize}

  \vspace{4pt}
  速度与位置更新：
  \begin{align*}
    \bm{v}_i &\leftarrow \omega \bm{v}_i + c_1 r_1 (\bm{p}_{\text{best},i} - \bm{x}_i) + c_2 r_2 (\bm{g}_{\text{best}} - \bm{x}_i) \\
    \bm{x}_i &\leftarrow \bm{x}_i + \bm{v}_i
  \end{align*}

  \begin{itemize}
    \item $\omega$：惯性权重（控制探索 vs 利用）
    \item $c_1, c_2$：认知 / 社会学习因子（典型值 $\approx 2$）
    \item 实数连续优化问题的主力算法之一
  \end{itemize}
\end{frame}

\subsection{禁忌搜索与蚁群算法}

\begin{frame}{禁忌搜索 (Tabu Search)}
  通过\textbf{禁忌列表 (Tabu List)}记忆近期访问过的解或操作，强制探索新区域。

  \vspace{6pt}
  \textbf{核心机制}：
  \begin{itemize}
    \item 即使邻域中没有更好的解，也移动到最优的“非禁忌”邻居
    \item 禁忌长度 (Tabu Tenure)：记录最近 $\ell$ 步禁止的移动
    \item 特赦准则 (Aspiration Criterion)：若禁忌移动能产生历史最佳解，可破例接受
    \item 长期记忆 (Long-term Memory)：强化好区域、惩罚探索过的区域（频率记忆）
  \end{itemize}

  特别擅长组合优化问题：TSP、调度、图着色等。
\end{frame}

\begin{frame}{蚁群算法 (Ant Colony Optimization, ACO)}
  模拟蚂蚁觅食的信息素机制：蚂蚁在走过的路径上释放信息素 (pheromone)，后续蚂蚁倾向于选择信息素浓度高的路径。

  \vspace{6pt}
  \textbf{TSP 上的 ACO}：
  \begin{itemize}
    \item 信息素 $\tau_{ij}$ 在边 $(i,j)$ 上积累
    \item 蚂蚁 $k$ 从城市 $i$ 选择下一城市 $j$ 的概率：
      \begin{equation*}
        p_{ij}^k = \frac{\tau_{ij}^\alpha \cdot \eta_{ij}^\beta}{\sum_{\ell \in N_i^k} \tau_{i\ell}^\alpha \cdot \eta_{i\ell}^\beta}
      \end{equation*}
      其中 $\eta_{ij} = 1/d_{ij}$（能见度 / 启发式信息）
    \item 信息素蒸发 (Evaporation)：$\tau_{ij} \leftarrow (1-\rho)\tau_{ij}$，避免过早收敛
    \item 精英蚂蚁加强最短路径上的信息素
  \end{itemize}

  广泛应用于路径规划、网络路由、调度等组合优化问题。
\end{frame}

\begin{frame}{启发式算法对比}
  \begin{table}
    \small
    \begin{tabular}{l|c|c|c}
      \toprule
      算法 & 搜索模式 & 记忆 & 典型问题 \\
      \midrule
      SA & 单点 + 概率跳跃 & 无 & 连续 / 组合 \\
      GA & 群体 + 交叉变异 & 隐性（种群） & 组合 / 连续 \\
      PSO & 群体 + 速度驱动 & $\bm{p}_{\text{best}}, \bm{g}_{\text{best}}$ & 连续 \\
      TS & 单点 + 禁忌引导 & 显性（列表） & 组合 \\
      ACO & 群体 + 信息素 & 分布式（信息素） & 路径 / 调度 \\
      \bottomrule
    \end{tabular}
  \end{table}

  \vspace{6pt}
  \textbf{No Free Lunch Theorem}：没有一个算法在所有问题上都最优。选择算法取决于问题结构、维度、约束类型。
\end{frame}

% PART 3: 凸优化
% ############################################################
\section{凸优化 (Convex Optimization)}

\subsection{凸集与凸函数}

\begin{frame}{凸集 (Convex Sets)}
  集合 $C \subseteq \R^n$ 为\textbf{凸集}，若对任意 $\bm{x}, \bm{y} \in C$ 和 $\theta \in [0,1]$：
  \begin{equation*}
    \theta \bm{x} + (1-\theta)\bm{y} \in C
  \end{equation*}

  \vspace{4pt}
  \begin{columns}[T]
    \column{0.48\textwidth}
    \textbf{常见凸集}：
    \begin{itemize}
      \item 超平面：$\{\bm{x} \mid \bm{a}^\top \bm{x} = b\}$
      \item 半空间：$\{\bm{x} \mid \bm{a}^\top \bm{x} \leq b\}$
      \item 多面体：$\{\bm{x} \mid \bm{Ax} \preceq \bm{b}\}$
      \item 椭球体
      \item 范数锥：$\{(\bm{x}, t) \mid \|\bm{x}\| \leq t\}$
    \end{itemize}
    \column{0.48\textwidth}
    \textbf{保凸运算}：
    \begin{itemize}
      \item 任意交集 $\bigcap_\alpha C_\alpha$
      \item 仿射映射 $f(C)$
      \item 透视函数 / 线性分式函数
    \end{itemize}
  \end{columns}
\end{frame}

\begin{frame}{凸函数 (Convex Functions)}
  函数 $f: \R^n \to \R$ 为\textbf{凸函数}，若其定义域是凸集且对任意 $\bm{x}, \bm{y}$ 和 $\theta \in [0,1]$：
  \begin{equation*}
    f(\theta \bm{x} + (1-\theta)\bm{y}) \leq \theta f(\bm{x}) + (1-\theta) f(\bm{y})
  \end{equation*}

  \vspace{6pt}
  \textbf{判别条件}：
  \begin{itemize}
    \item 一阶条件（可微时）：$f(\bm{y}) \geq f(\bm{x}) + \nabla f(\bm{x})^\top (\bm{y}-\bm{x})$
    \item 二阶条件（二阶可微时）：$\nabla^2 f(\bm{x}) \succcurlyeq \bm{0}$（Hessian 半正定）
    \item 上镜图 (Epigraph)：$\epi(f) = \{(\bm{x}, t) \mid f(\bm{x}) \leq t\}$ 是凸集 $\Longleftrightarrow$ $f$ 是凸函数
  \end{itemize}

  \vspace{4pt}
  \textbf{强凸函数}：$f(\bm{y}) \geq f(\bm{x}) + \nabla f(\bm{x})^\top (\bm{y}-\bm{x}) + \frac{\mu}{2}\|\bm{y}-\bm{x}\|^2$，$\mu > 0$ 唯一极小点。
\end{frame}

\begin{frame}{凸优化问题}
  标准凸优化问题：
  \begin{align*}
    \min_{\bm{x}} \quad & f_0(\bm{x}) \\
    \text{s.t.} \quad   & f_i(\bm{x}) \leq 0,\quad i=1,\dots,m \\
                        & \bm{Ax} = \bm{b}
  \end{align*}
  其中 $f_0, f_1, \dots, f_m$ 均为凸函数。

  \vspace{6pt}
  \begin{block}{关键性质}
    任何局部最优解 = 全局最优解。KKT 条件既是必要也是充分的。
  \end{block}

  \vspace{4pt}
  \textbf{常见凸优化问题}：
  \begin{itemize}
    \item LP, QP (二次规划), QCQP (二次约束二次规划)
    \item SOCP (二阶锥规划), SDP (半定规划)
    \item 几何规划 (Geometric Programming)
  \end{itemize}
\end{frame}

\begin{frame}{凸对偶理论}
  Lagrange 对偶函数：
  \begin{equation*}
    g(\bm{\lambda}, \bm{\nu}) = \inf_{\bm{x}} \mathcal{L}(\bm{x}, \bm{\lambda}, \bm{\nu})
  \end{equation*}
  其中 $\mathcal{L}(\bm{x}, \bm{\lambda}, \bm{\nu}) = f_0(\bm{x}) + \sum_i \lambda_i f_i(\bm{x}) + \sum_j \nu_j (\bm{a}_j^\top \bm{x} - b_j)$。

  \vspace{6pt}
  \begin{block}{Slater 条件 (Slater's Condition)}
    若存在严格可行点（所有不等式约束严格成立），则强对偶成立：
    \begin{equation*}
      \text{Primal 最优值} = \text{Dual 最优值}
    \end{equation*}
  \end{block}
  对偶问题始终是凹的（极大化），且自动给出原始问题的下界。
\end{frame}

\subsection{凸优化算法}

\begin{frame}{次梯度方法 (Subgradient Method)}
  对于不可微的凸函数 $f$：
  \begin{equation*}
    \bm{x}_{k+1} = \bm{x}_k - \alpha_k \bm{g}_k,\quad \bm{g}_k \in \partial f(\bm{x}_k)
  \end{equation*}
  其中 $\partial f(\bm{x})$ = 次梯度集合。

  \begin{itemize}
    \item 步长满足 $\alpha_k \to 0,\ \sum \alpha_k = \infty$ 时收敛
    \item 收敛速度 $\mathcal{O}(1/\sqrt{k})$
    \item $\ell_1$ 正则化、SVM 等场景常用
  \end{itemize}

  \vspace{6pt}
  \begin{alertblock}{注意}
    次梯度不一定下降方向，需跟踪迄今最佳点（而非最后一步）。
  \end{alertblock}
\end{frame}

\begin{frame}{近端梯度法 (Proximal Gradient Method)}
  目标：$\min f(\bm{x}) + g(\bm{x})$，$f$ 光滑凸，$g$ 可能不可微但 prox 易算。

  \vspace{4pt}
  \textbf{近端算子 (Proximal Operator)}：
  \begin{equation*}
    \prox_{\eta g}(\bm{v}) = \argmin_{\bm{x}} \left\{ g(\bm{x}) + \frac{1}{2\eta}\|\bm{x} - \bm{v}\|^2 \right\}
  \end{equation*}

  \vspace{4pt}
  \textbf{ISTA / FISTA}（迭代 / 快速收缩阈值算法）：
  \begin{align*}
    \bm{z}_k &= \bm{x}_k - \eta \nabla f(\bm{x}_k) \\
    \bm{x}_{k+1} &= \prox_{\eta g}(\bm{z}_k)
  \end{align*}

  \begin{itemize}
    \item 适用于 $\ell_1$ 正则化 (LASSO)、低秩矩阵恢复等
    \item FISTA 加速至 $\mathcal{O}(1/k^2)$
  \end{itemize}
\end{frame}

\begin{frame}{Nesterov 加速梯度 (Accelerated Gradient)}
  对于光滑凸函数，经典梯度下降 $\mathcal{O}(1/k)$，而 Nesterov 加速可达 $\mathcal{O}(1/k^2)$。

  \vspace{4pt}
  \textbf{Nesterov 加速梯度 (NAG)}：
  \begin{align*}
    \bm{y}_k &= \bm{x}_k + \beta_k(\bm{x}_k - \bm{x}_{k-1}) \\
    \bm{x}_{k+1} &= \bm{y}_k - \alpha_k \nabla f(\bm{y}_k)
  \end{align*}
  其中 $\beta_k = \frac{k-1}{k+2}$（经典方案）。

  \begin{itemize}
    \item "Look ahead" 策略：不在当前位置求梯度，而在外推位置
    \item 对强凸函数可达线性收敛速率 $\mathcal{O}\big((1 - \sqrt{\mu/L})^k\big)$
    \item Nesterov 动量是 SGD with Momentum 和 Adam 的前身
  \end{itemize}
\end{frame}

\begin{frame}{ADMM (Alternating Direction Method of Multipliers)}
  适用于结构分离的凸优化：
  \begin{align*}
    \min_{\bm{x}, \bm{z}} \quad & f(\bm{x}) + g(\bm{z}) \\
    \text{s.t.} \quad           & \bm{Ax} + \bm{Bz} = \bm{c}
  \end{align*}

  \vspace{4pt}
  增广 Lagrange：$\mathcal{L}_\rho(\bm{x}, \bm{z}, \bm{y}) = f(\bm{x}) + g(\bm{z}) + \bm{y}^\top(\bm{Ax} + \bm{Bz} - \bm{c}) + \frac{\rho}{2}\|\bm{Ax} + \bm{Bz} - \bm{c}\|^2$

  \vspace{4pt}
  ADMM 迭代：
  \begin{align*}
    \bm{x}_{k+1} &= \argmin_{\bm{x}} \mathcal{L}_\rho(\bm{x}, \bm{z}_k, \bm{y}_k) \\
    \bm{z}_{k+1} &= \argmin_{\bm{z}} \mathcal{L}_\rho(\bm{x}_{k+1}, \bm{z}, \bm{y}_k) \\
    \bm{y}_{k+1} &= \bm{y}_k + \rho(\bm{Ax}_{k+1} + \bm{Bz}_{k+1} - \bm{c})
  \end{align*}

  分布式优化、图像处理、统计学习等领域广泛应用。
\end{frame}

\begin{frame}{半定规划 (SDP) 简介}
  SDP 是 LP 向矩阵空间的推广：
  \begin{align*}
    \min_{\bm{X}} \quad & \langle \bm{C}, \bm{X} \rangle = \tr(\bm{C}^\top \bm{X}) \\
    \text{s.t.} \quad   & \langle \bm{A}_i, \bm{X} \rangle = b_i,\quad i=1,\dots,m \\
                        & \bm{X} \succcurlyeq \bm{0} \quad (\text{半正定约束})
  \end{align*}

  \begin{itemize}
    \item 内点法 (Primal-Dual IPM) 可在多项式时间内求解
    \item 应用：矩阵补全、鲁棒 PCA、组合优化的 SDP 松弛
    \item Max-Cut 的 Goemans–Williamson 0.878 近似比即基于 SDP
  \end{itemize}

  LP $\subset$ SOCP $\subset$ SDP，逐层增加表达能力。
\end{frame}

\begin{frame}{凸优化求解器概览}
  \begin{columns}[T]
    \column{0.48\textwidth}
    \textbf{通用求解器}：
    \begin{itemize}
      \item CVXPY (Python)
      \item CVX (MATLAB)
      \item Convex.jl (Julia)
    \end{itemize}
    \column{0.48\textwidth}
    \textbf{专用求解器}：
    \begin{itemize}
      \item ECOS, SCS — SOCP, SDP
      \item OSQP — QP
      \item MOSEK — 商业级 LP/QP/SOCP/SDP
      \item Gurobi / CPLEX — LP, MILP, QP
    \end{itemize}
  \end{columns}

  \vspace{8pt}
  \textbf{建模原则}：
  \begin{itemize}
    \item 若能建模为 LP → 用 LP 求解器（最快）
    \item 若能建模为 QP → 用 QP 求解器
    \item 若能建模为 SOCP/SDP → 用相应锥规划求解器
    \item 只有非凸时才考虑通用 NLP 求解器 (IPOPT, SNOPT 等)
  \end{itemize}
\end{frame}

% ############################################################
% PART 4: 运筹优化算法在机器学习和大模型中的应用
% ############################################################
\section{运筹优化在 ML 与大模型中的应用}

\subsection{ML 中的优化问题}

\begin{frame}{机器学习 = 优化}
  机器学习的核心是优化经验风险：
  \begin{equation*}
    \min_{\bm{\theta}} \frac{1}{N} \sum_{i=1}^N \ell(\bm{x}_i, y_i;\ \bm{\theta}) + \lambda R(\bm{\theta})
  \end{equation*}

  \begin{itemize}
    \item $\ell(\cdot)$：损失函数（交叉熵、MSE、Hinge 等）
    \item $R(\cdot)$：正则化项（$\ell_2$, $\ell_1$, dropout 等）
    \item $N$：样本量（现代 ML 中常为百万~十亿级）
  \end{itemize}

  \vspace{6pt}
  \textbf{ML 优化的特点}：
  \begin{itemize}
    \item $N$ 极大 $\to$ 随机梯度方法
    \item $n$（参数维度）极大 $\to$ 一阶方法为主
    \item 非凸但实践中可找到不错的局部极小点
    \item 对泛化误差的关心 $\gg$ 对训练误差的关心
  \end{itemize}
\end{frame}

\begin{frame}{随机梯度下降 (SGD)}
  每次用一个小批量 (mini-batch) $\mathcal{B}$ 估计梯度：
  \begin{equation*}
    \bm{\theta}_{t+1} = \bm{\theta}_t - \eta_t \cdot \frac{1}{|\mathcal{B}|} \sum_{i \in \mathcal{B}} \nabla_\theta \ell(\bm{x}_i, y_i;\ \bm{\theta}_t)
  \end{equation*}

  \begin{itemize}
    \item 无偏估计：$\E[\nabla f_{\mathcal{B}}(\bm{\theta})] = \nabla f(\bm{\theta})$
    \item 收敛速度 $\mathcal{O}(1/\sqrt{T})$（非凸），$\mathcal{O}(1/T)$（强凸）
    \item 关键技巧：学习率调度、batch size 选择
  \end{itemize}

  \begin{alertblock}{瓶颈}
    纯 SGD 收敛慢，对 ill-conditioned 和梯度噪声敏感。
  \end{alertblock}
\end{frame}

\begin{frame}{动量法 (Momentum) 与 Nesterov}
  \textbf{SGD with Momentum (Polyak)}：
  \begin{align*}
    \bm{v}_{t} &= \beta \bm{v}_{t-1} + \eta \nabla f(\bm{\theta}_t) \\
    \bm{\theta}_{t+1} &= \bm{\theta}_t - \bm{v}_t
  \end{align*}

  \textbf{Nesterov Accelerated Gradient (NAG)}：
  \begin{align*}
    \bm{v}_{t} &= \beta \bm{v}_{t-1} + \eta \nabla f(\bm{\theta}_t - \beta \bm{v}_{t-1}) \\
    \bm{\theta}_{t+1} &= \bm{\theta}_t - \bm{v}_t
  \end{align*}

  \begin{itemize}
    \item 动量累积历史梯度方向，加速收敛、减少震荡
    \item NAG 先根据动量外推，再计算梯度（"Look-ahead"）
    \item 典型 $\beta = 0.9$
  \end{itemize}
\end{frame}

\begin{frame}{Adam (Adaptive Moment Estimation)}
  结合动量与自适应学习率：
  \begin{align*}
    \bm{m}_t &= \beta_1 \bm{m}_{t-1} + (1-\beta_1) \bm{g}_t \\
    \bm{v}_t &= \beta_2 \bm{v}_{t-1} + (1-\beta_2) \bm{g}_t^2 \\
    \hat{\bm{m}}_t &= \bm{m}_t / (1-\beta_1^t),\quad \hat{\bm{v}}_t = \bm{v}_t / (1-\beta_2^t) \\
    \bm{\theta}_{t+1} &= \bm{\theta}_t - \eta \cdot \frac{\hat{\bm{m}}_t}{\sqrt{\hat{\bm{v}}_t} + \epsilon}
  \end{align*}

  \begin{itemize}
    \item $\hat{\bm{m}}_t$：一阶动量（指数移动平均）
    \item $\hat{\bm{v}}_t$：二阶动量（RMS 缩放，类似对角的二阶方法近似）
    \item 默认参数：$\beta_1=0.9$, $\beta_2=0.999$, $\epsilon=10^{-8}$
    \item 变体：AdamW（解耦 weight decay）、AdaBelief、LAMB
  \end{itemize}
\end{frame}

\begin{frame}{学习率调度 (Learning Rate Schedule)}
  学习率 $\eta$ 对收敛至关重要：

  \begin{itemize}
    \item \textbf{Step decay}：每隔固定 epoch 乘以衰减因子
    \item \textbf{Cosine annealing}：
      \begin{equation*}
        \eta_t = \eta_{\min} + \frac12(\eta_{\max} - \eta_{\min})\left(1 + \cos\frac{t\pi}{T}\right)
      \end{equation*}
    \item \textbf{Warmup}：前几个 epoch 线性增大 $\eta$（Transformer 训练关键）
    \item \textbf{Cyclic / OneCycle}：周期性调整 LR，加速收敛
  \end{itemize}

  \vspace{4pt}
  \begin{block}{Warmup 的必要性}
    Transformer 训练初期，Warmup 稳定训练、防止梯度爆炸；没有 warmup 时 Adam 在初始阶段方差过大。
  \end{block}
\end{frame}

\begin{frame}{二阶方法与深度学习}
  深度学习主要用一阶方法，但二阶方法有吸引力：
  \begin{itemize}
    \item \textbf{Natural Gradient} / \textbf{K-FAC}：用 Fisher 信息矩阵近似 Hessian
    \item \textbf{Shampoo}：逐层预处理，Google 大规模训练使用
    \item \textbf{AdaHessian}：用 Hutchinson 方法随机估计 Hessian 对角线
    \item \textbf{Newton Sketch}：随机投影加速 Newton 法
  \end{itemize}

  \vspace{6pt}
  挑战：Hessian 维度 $n \times n$ 太大（数十亿参数 $\to$ 不可能存储），且深度网络 Hessian 谱复杂（负特征值、大条件数）。
\end{frame}

\subsection{大规模分布式优化}

\begin{frame}{分布式训练中的优化}
  大模型（> 100B 参数）需要跨数千 GPU 的分布式训练。

  \vspace{4pt}
  \textbf{数据并行 (Data Parallel)} + \textbf{All-Reduce}：
  \begin{itemize}
    \item 每 GPU 持有完整模型副本，处理不同数据 mini-batch
    \item 求完梯度后 All-Reduce 求平均、更新
  \end{itemize}

  \vspace{4pt}
  \textbf{模型并行 (Model Parallel)}：
  \begin{itemize}
    \item Tensor Parallelism：矩阵乘法在多个 GPU 上分片
    \item Pipeline Parallelism：不同层在不同 GPU 上
  \end{itemize}

  \vspace{4pt}
  \textbf{关键优化技术}：
  \begin{itemize}
    \item \textbf{Gradient Accumulation}：多步累积梯度后更新（模拟大 batch）
    \item \textbf{混合精度 (FP16/BF16)}：降低通信和计算开销
    \item \textbf{ZeRO} (Zero Redundancy Optimizer)：分片优化器状态
  \end{itemize}
\end{frame}

\begin{frame}{ZeRO 优化器 (DeepSpeed)}
  Microsoft DeepSpeed 的 ZeRO 将优化器状态分区以节省 GPU 显存：

  \begin{table}
    \small
    \begin{tabular}{l|p{5.5cm}|p{5cm}}
      \toprule
      阶段 & 分片内容 & 显存节省 \\
      \midrule
      ZeRO-1 & 优化器状态 (Adam $m, v$) & 4$\times$ \\
      ZeRO-2 & 优化器状态 + 梯度 & 8$\times$ \\
      ZeRO-3 & 优化器状态 + 梯度 + 参数 & 线性随 GPU 数扩展 \\
      \bottomrule
    \end{tabular}
  \end{table}

  \vspace{4pt}
  这使 175B 参数的 GPT-3 级模型可以在数百个 GPU 上训练，而不需要数千个。
\end{frame}

\subsection{LLM 微调中的优化}

\begin{frame}{LLM 微调：优化视角}
  全参数微调 (Full Fine-tuning) 对千亿参数模型代价极高：
  \begin{itemize}
    \item 需要存储所有参数梯度
    \item 需要 16-32 倍模型参数量的显存（优化器状态 + 梯度 + 激活值）
  \end{itemize}

  \vspace{6pt}
  \textbf{参数高效微调 (PEFT)} 大幅减少可训练参数：
  \begin{itemize}
    \item \textbf{LoRA}：低秩矩阵分解，仅训练低秩增量 $\Delta W = BA$
    \item \textbf{Adapter}：在层间插入小型瓶颈网络
    \item \textbf{Prefix / Prompt Tuning}：仅优化虚拟 token 嵌入
  \end{itemize}

  本质上都是\textbf{降维优化}——约束优化在一个低维子流形上进行。
\end{frame}

\begin{frame}{RLHF 中的优化}
  RLHF (Reinforcement Learning from Human Feedback) 的优化流程：
  \begin{enumerate}
    \item SFT (Supervised Fine-Tuning)：标准监督学习
    \item 训练 Reward Model：Bradley-Terry 偏好模型
    \item PPO (Proximal Policy Optimization)：用 Reward Model 优化策略
      \begin{equation*}
        \max_\theta \E\left[ r(x,y) - \beta \cdot \text{KL}(\pi_\theta \parallel \pi_{\text{ref}}) \right]
      \end{equation*}
  \end{enumerate}

  \vspace{4pt}
  替代方案：\textbf{DPO (Direct Preference Optimization)}——将 RL 优化转化为一个直接的分类损失，避免了显式的 reward 建模和 PPO 训练。

  \begin{equation*}
    \mathcal{L}_{\text{DPO}}(\theta) = -\E \left[ \log \sigma \left( \beta \log \frac{\pi_\theta(y_w|x)}{\pi_{\text{ref}}(y_w|x)} - \beta \log \frac{\pi_\theta(y_l|x)}{\pi_{\text{ref}}(y_l|x)} \right) \right]
  \end{equation*}
\end{frame}

\begin{frame}{量化与推理优化}
  大模型推理部署的优化技术：

  \begin{itemize}
    \item \textbf{权重量化 (GPTQ, AWQ)}：将 FP16 权重压缩到 4-bit 或 8-bit
      \begin{itemize}
        \item 基于最优脑外科 (OBS) 的逐层量化
        \item 目标：$\min_{\hat{W}} \|WX - \hat{W}X\|_F^2$
      \end{itemize}
    \item \textbf{KV Cache 优化}：量化缓存的 Key-Value 对
    \item \textbf{投机解码 (Speculative Decoding)}：用小模型生成候选 token，大模型验证
    \item \textbf{FlashAttention}：IO-aware 的精确注意力计算
  \end{itemize}

  \vspace{4pt}
  这些都是经典的\textbf{近似优化}问题：在精度、速度、内存之间寻找 Pareto 最优。
\end{frame}

\begin{frame}{大模型训练中的优化挑战}
  LLM 训练的优化难点：
  \begin{itemize}
    \item \textbf{损失尖峰 (Loss Spike)}：训练不稳定，需要回滚和重新开始
    \item \textbf{数值精度}：FP16/BF16 下的梯度 underflow / overflow
    \item \textbf{梯度裁剪 (Gradient Clipping)}：$\|\bm{g}\| > \tau$ 时缩放 $\bm{g} \leftarrow \tau \cdot \bm{g} / \|\bm{g}\|$
    \item \textbf{Z-loss / Logit Regularization}：防止 logit 漂移
    \item \textbf{嵌入层归一化}：output embedding 的 LayerNorm 稳定训练
  \end{itemize}

  \vspace{6pt}
  \textbf{Chinchilla Scaling Law}：在计算预算 $C$ 固定时，模型参数 $N$ 与训练 token 数 $D$ 的最优配置：
  \begin{equation*}
    N_{\text{opt}} \propto C^{0.5},\quad D_{\text{opt}} \propto C^{0.5}
  \end{equation*}
\end{frame}

\begin{frame}{运筹优化与 ML 的交叉前沿}
  \begin{itemize}
    \item \textbf{Decision-Focused Learning}：将 ML 预测嵌入优化模型，端到端训练
      \begin{itemize}
        \item 例：预测需求 $\to$ 库存优化，梯度需穿过 LP/IP 层
      \end{itemize}
    \item \textbf{Learning to Optimize}：用 RL 学习分支定界启发式、切平面选择等
    \item \textbf{组合优化的神经求解器}：Pointer Network / Transformer 解 TSP, VRP
    \item \textbf{混合整数规划 + ML}：用 GNN 加速 MILP 求解
    \item \textbf{Convex ML}：凸损失 + 凸正则项的模型（如 Elastic Net）
    \item \textbf{Bilevel Optimization}：元学习、超参数优化、对抗训练
  \end{itemize}
\end{frame}

% ############################################################
% PART 5: GPU/CUDA 并行计算与 ML 优化加速
% ############################################################
\section{GPU / CUDA 并行计算与 ML 优化加速}

\begin{frame}{为什么 GPU 对优化算法至关重要？}
  现代 ML 优化的核心计算是\textbf{大规模矩阵运算}和\textbf{梯度计算}——天然适合 GPU 的 SIMT 架构。

  \vspace{6pt}
  \begin{columns}[T]
    \column{0.48\textwidth}
    \textbf{CPU (少核大算力)}：
    \begin{itemize}
      \item 8--64 个复杂核心
      \item 擅长分支、串行逻辑
      \item 大缓存、乱序执行
    \end{itemize}
    \column{0.48\textwidth}
    \textbf{GPU (多核轻量)}：
    \begin{itemize}
      \item 数千个简单核心
      \item 擅长数据并行 (SIMT)
      \item 高内存带宽 ($\sim$2 TB/s)
    \end{itemize}
  \end{columns}

  \vspace{8pt}
  \begin{block}{NVIDIA CUDA 编程模型}
    \begin{itemize}
      \item Thread → Warp (32 threads) → Block → Grid
      \item 内存层次：Global → L2 → Shared Memory → Registers
      \item Tensor Cores (Volta+)：单周期 $4\times4$ 矩阵乘加
    \end{itemize}
  \end{block}

  \textcolor{blue}{A100 80GB: 6912 CUDA cores + 432 Tensor Cores, 312 TFLOPS (FP16)}
\end{frame}

\begin{frame}{GPU 上的优化算法：一阶方法的天然优势}
  梯度下降族算法 (SGD, Adam, AdamW) 的每次迭代：
  \begin{enumerate}
    \item Forward pass：矩阵乘法 $\bm{y} = \bm{Wx}$
    \item Loss 计算：标量，计算量小
    \item Backward pass：矩阵乘法 $\nabla_W \mathcal{L} = (\nabla_y \mathcal{L}) \bm{x}^\top$
    \item 参数更新：逐元素操作
  \end{enumerate}

  \vspace{6pt}
  \textbf{关键观察}：第 1、3 步是密集矩阵乘法，GPU 加速 $50\times$--$200\times$。

  \begin{itemize}
    \item cuBLAS：高度优化的 GEMM (通用矩阵乘)
    \item cuDNN：卷积、池化、归一化的 GPU kernel
    \item 二阶方法 (Newton, K-FAC) 的 Hessian 构造在 GPU 上仍昂贵
  \end{itemize}

  \begin{alertblock}{GPU 瓶颈}
    不是计算，是\textbf{显存带宽}和\textbf{通信}。一个 175B 模型需要 $\sim$350 GB 仅存参数 (FP16)，加上优化器状态轻松破 TB。
  \end{alertblock}
\end{frame}

\begin{frame}{混合精度训练 (Mixed Precision)}
  \textbf{动机}：FP32 精度在很多梯度更新中冗余，FP16 可翻倍吞吐、减半显存。

  \vspace{4pt}
  \textbf{标准混合精度流程 (NVIDIA APEX / PyTorch AMP)}：
  \begin{enumerate}
    \item 权重以 FP32 主副本存储
    \item Forward/Backward 用 FP16 计算（速度 $2\times$）
    \item \textbf{Loss Scaling}：梯度太小时乘以常数 $S$ 防止 underflow
    \item 梯度回转为 FP32，更新主副本权重
  \end{enumerate}

  \begin{block}{BF16 (Brain Float 16)}
    \begin{itemize}
      \item 与 FP32 相同指数位 (8 bit)，仅尾数截断 (7 bit)
      \item 动态范围 = FP32，无需 loss scaling
      \item NVIDIA A100 / H100 原生支持，训练更稳定
    \end{itemize}
  \end{block}

  Tensor Cores 在 FP16/BF16 下吞吐可达 FP32 的 8--16$\times$。
\end{frame}

\begin{frame}{梯度累积 \& 大批量训练}
  GPU 显存限制单次能处理的 batch size，但大批量训练有助于收敛稳定性和扩展性。

  \vspace{6pt}
  \textbf{梯度累积 (Gradient Accumulation)}：
  \begin{itemize}
    \item 执行 $k$ 次 micro-batch 前向+反向，梯度累加但不更新
    \item 第 $k$ 次后统一更新，等价于 batch size $\times k$
    \item 显存需求 = micro-batch size，训练时间 $\approx$ 不变（梯度同步开销略增）
  \end{itemize}

  \vspace{4pt}
  \textbf{大批量训练的挑战}：
  \begin{itemize}
    \item 线性缩放规则 (Linear Scaling Rule, Goyal et al. 2017)：\\
      $\eta_{\text{new}} = \eta_{\text{base}} \times (B_{\text{new}} / B_{\text{base}})$，需配合 warmup
    \item LARS (Layer-wise Adaptive Rate Scaling)：逐层自适应学习率
    \item LAMB：Adam 的大批量变体，BERT 预训练标配
  \end{itemize}
\end{frame}

\begin{frame}{分布式 GPU 训练架构}
  \begin{columns}[T]
    \column{0.55\textwidth}
    \textbf{数据并行 (Data Parallel)}：
    \begin{itemize}
      \item 每 GPU 持有完整模型
      \item 各 GPU 吃不同 mini-batch
      \item All-Reduce 求梯度平均
      \item PyTorch DDP, Horovod
    \end{itemize}

    \vspace{4pt}
    \textbf{模型并行}：
    \begin{itemize}
      \item Tensor Parallel (Megatron-LM)
      \item Pipeline Parallel (GPipe)
      \item 3D Parallel = DP + TP + PP
    \end{itemize}
    \column{0.45\textwidth}
    \textbf{NCCL 通信原语}：
    \begin{itemize}
      \item All-Reduce：梯度平均
      \item All-Gather：收集分片
      \item Reduce-Scatter：分片归约
      \item 带宽：$\sim$600 GB/s (NVLink)
    \end{itemize}

    \vspace{4pt}
    \textbf{FSDP (Fully Sharded DP)}：\\
    参数+梯度+优化器全分片，ZeRO-3 的 PyTorch 原生实现。
  \end{columns}
\end{frame}

\begin{frame}{显存优化策略}
  训练大模型的显存构成（以 1B 参数 + Adam + FP16 为例）：

  \begin{table}
    \small
    \begin{tabular}{l|r|r}
      \toprule
      组成部分 & 估算大小 & 比例 \\
      \midrule
      模型参数 (FP16) & 2 GB & $\sim$20\% \\
      Adam $m$ (FP32) & 4 GB & $\sim$40\% \\
      Adam $v$ (FP32) & 4 GB & $\sim$40\% \\
      梯度 (FP16) & 2 GB & $\sim$20\% \\
      激活值 (Activations) & 可变 & 额外 \\
      \bottomrule
    \end{tabular}
  \end{table}

  \vspace{4pt}
  \textbf{核心优化技术}：
  \begin{itemize}
    \item \textbf{Gradient Checkpointing}：不存中间激活，反向时重新计算（显存 $\searrow$、计算 $\nearrow$）
    \item \textbf{Activation Offloading}：激活值溢出到 CPU 内存
    \item \textbf{FlashAttention}：IO-aware 精确注意力，分块计算 + 在线 softmax，显存从 $\mathcal{O}(N^2)$ 降到 $\mathcal{O}(N)$
    \item \textbf{Kernel Fusion}：融合多个 CUDA kernel 减少访存（如 Bias + GeLU + Dropout）
  \end{itemize}
\end{frame}

\begin{frame}{GPU 优化实战 Checklist}
  \begin{enumerate}
    \item \textbf{混合精度}：启用 AMP (FP16) 或 BF16，确认 Tensor Core 利用率
    \item \textbf{数据加载}：多进程 DataLoader，prefetch，pin\_memory
    \item \textbf{梯度累积}：用小 batch + 累积模拟大 batch
    \item \textbf{分布式}：单机多卡用 DDP，多机用 FSDP / DeepSpeed
    \item \textbf{编译优化}：torch.compile / TensorRT，XLA，kernel fusion
    \item \textbf{显存管理}：gradient checkpointing，FlashAttention，fused optimizer
    \item \textbf{监控工具}：nvidia-smi, Nsight Systems (timeline), Nsight Compute (kernel)
    \item \textbf{batch size 调优}：找到 GPU 利用率拐点 (通常 $> 80\%$ 为佳)
  \end{enumerate}

  \vspace{6pt}
  \begin{center}
    \textcolor{blue}{\textbf{核心原则：减少 GPU 空闲等待 —— 要么在算，要么在传数据，不要闲下来。}}
  \end{center}
\end{frame}

\begin{frame}{GPU 加速的运筹优化应用}
  \begin{itemize}
    \item \textbf{大规模 LP 求解}：GPU 加速内点法的线性系统求解（cuSOLVER）
    \item \textbf{分支定界并行}：MIP 的 B\&B 树在 GPU 上并行展开
    \item \textbf{组合优化的神经求解器}：Transformer/GNN 在 GPU 上做 TSP/VRP 端到端推理
    \item \textbf{蒙特卡洛仿真}：GPU 加速随机优化（如随机规划的场景生成与求解）
    \item \textbf{RL for OR}：GPU 并行仿真环境，PPO 训练调度/路径策略
    \item \textbf{SDP 求解}：GPU 加速 SDP 内点法中的 Cholesky 分解和矩阵乘法
  \end{itemize}

  \vspace{8pt}
  \begin{block}{趋势}
    NVIDIA cuOpt：专为运筹优化设计的 GPU 加速库，支持 VRP、调度等问题，速度可达传统求解器的 $100\times$。
  \end{block}
\end{frame}


\begin{frame}{总结}
  \begin{enumerate}
    \item \textbf{线性规划}是优化理论的基石，单纯形法与内点法互补
    \item \textbf{非线性规划}：KKT 条件统一处理约束问题，一阶/二阶方法各有适用场景
    \item \textbf{凸优化}：局部即全局，加速梯度、ADMM、近端方法解决大规模问题
    \item \textbf{ML 中的优化}：SGD → Momentum → Adam，一阶方法主导
    \item \textbf{LLM 中的优化}：分布式训练（ZeRO）、PEFT（LoRA）、量化、RLHF
    \item \textbf{交叉前沿}：ML 与 OR 双向赋能，端到端优化决策
    \item \textbf{GPU 加速}：CUDA、混合精度、分布式训练、FlashAttention、cuOpt
  \end{enumerate}

  \vspace{12pt}
  \begin{center}
    \Large \textbf{谢谢！}
  \end{center}
\end{frame}

\end{document}
